\documentclass[12pt]{article}
\usepackage[margin=1in]{geometry}

\begin{document}

\newpage
\section*{Novelty Statement}

This paper makes an original methodological contribution to tourism economics by applying a synthetic counterfactual framework to evaluate the labor market effects of a negative tourism supply shock in a small island economy where ideal control groups and uniform data are not available. The approach addresses a common challenge in tourism research: identifying credible comparisons for destinations with unique economic structures and limited comparable units. Using monthly passenger flows, labor market indicators, and external comparison units, the study constructs a synthetic control to estimate the counterfactual path of Iceland’s tourism-linked employment following the bankruptcy of WOW Air. The methodology is further strengthened through placebo tests and sensitivity analyses to verify robustness and applicability.

The paper contributes to the intellectual agenda of tourism resilience by showing how transportation infrastructure shocks can generate persistent labor losses beyond the directly affected firm. The findings highlight the broader vulnerability of tourism-dependent economies and provide a replicable method for evaluating similar shocks in destinations where traditional identification strategies are constrained by data limitations and structural uniqueness.

\end{document}